\noindent A companion work developed a generative relational epistemology, on which a concept becomes observable through the relations among the several symbolic systems that render it, and knowing takes the form of a generative observation that is historically and dialectically produced. That work observed, near its end, that a person is the bearer of a symbolic system, so that the meeting of two persons is a crossing of symbolic worlds. This paper takes up that observation in the domain of intimacy, where the crossing is deepest and its stakes are highest. It asks what relational understanding, so conceived, does in an intimate bond, whether it is good, and how the capacities it requires might be cultivated. The paper builds two cases and then a practice. The functional case sets the potential work of relational understanding, in the generation of trust, the recognition of shared value, the easing of a certain class of conflict, and the quality of a shared life, beside the established theories and findings that treat these same matters, and it locates the distinctive contribution in two characters of the account, the relational and the generative, showing what each supplies that the received theories leave unsupplied. The normative case then argues that a functional defense is not sufficient, since an efficient possession of the other can counterfeit the functions of understanding, and it grounds a normative distinction in ethical legitimacy, observed across several ethical traditions and not rested on a single one. Consequentialist, Kantian, care-ethical, and Aristotelian renderings each supply an aspect that the others leave, and non-possession, with generative justice, stands among them as one aspect, so that the normative ground is drawn in the same relational manner as the epistemology, an instance of the very understanding the paper concerns. Consistent with the dialectical and materialist conception of truth the companion work develops, the paper does not claim the correctness of a given understanding, and holds instead that understanding, like truth, is generated in historical and contradictory movement, while what can be asked of it is that it be ethically legitimate, a legitimacy whose own criteria are themselves historically generated and set as no final guarantee. The paper then states the special problem of intimacy, that the bond is a site where these capacities are wagered and not rehearsed, since the value of what is given may return only after a long passage of time and cannot be tried again, so that the field in which a capacity is cultivated and the field in which it is exercised come apart. From this it proposes a model, a set of capacities and not a procedure, divided into those that admit near practice and those that can only be cultivated as a standing disposition, among the latter the throwing of a generative thing into the other's own generative repertoire, the non-intervention that lets a shared dynamic unfold, and the preservation of a relational experience against the day its value is recognized. Non-possession runs throughout as the condition without which the capacities decline into a refined extraction. The paper then turns dialectically upon its own argument, and holds, on a ground drawn from the psychoanalytic account of desire and from its own account of generation, that a full crossing of the other's world would end the bond it was meant to perfect, since the other's unsymbolizable remainder is the core around which desire moves and the reservoir from which the bond's further generation is drawn, so that understanding rightly aimed seeks the other and honors, in the same motion, what in the other it will not reach. The paper offers no final synthesis, and its form keeps the openness its account of understanding requires.
